% ══════════════════════════════════════════════════════════════════════
% RELATÓRIO DE AVALIAÇÃO DO ASSISTENTE DO RELATÓRIO GESTÃO HOSPITALAR NO SUS-SP
% ══════════════════════════════════════════════════════════════════════
% Cópia de estilo/modelo_relatorio.tex (o kit do portfólio). O bloco entre
% "PREÂMBULO DO KIT" e "FIM DO PREÂMBULO DO KIT" é o do kit, sem edição.
%
% Todo número vem de uma macro \Num... de numeros.tex, gerado por
% scripts/06_gerar_documentos.py a partir de resultados/numeros.json, e toda
% tabela vem de um arquivo de linhas de tabelas/, escrito por
% scripts/04_analisar.py e lido com \linhasde. As figuras vêm de ../figuras/,
% escritas por scripts/05_figuras.py. Nenhum número é digitado aqui.
%
% Compilar: uv run python scripts/reproduzir.py roda os scripts 04, 05 e 06,
% compila e confere. O pdflatex roda dentro da pasta relatorio/ da raiz (a
% do projeto ou a de RAG_SUS_RAIZ), de onde lê numeros.tex, tabelas/ e
% ../figuras/; este .tex nunca é copiado. À mão, de dentro de relatorio/:
%   pdflatex -interaction=nonstopmode -halt-on-error relatorio.tex   (duas vezes)
% Depois, scripts/conferir_numeros.py confere que nenhum número foi digitado.

% ══════════════════════════════════════════════════════════════════════
% PREÂMBULO DO KIT: daqui até "FIM DO PREÂMBULO DO KIT", não edite
% ══════════════════════════════════════════════════════════════════════
\documentclass[11pt,a4paper]{article}

\usepackage[T1]{fontenc}
\usepackage[utf8]{inputenc}
\usepackage[brazilian]{babel}

% A mancha: 15 cm de linha, cerca de 3,2 alfabetos de Libertine 11pt. As
% figuras NÃO vivem nesta medida: ver \figwidth, mais abaixo.
\usepackage[top=2cm, bottom=2cm, left=3cm, right=3cm]{geometry}
\usepackage{graphicx}
\usepackage{subcaption}
\usepackage{booktabs}
\usepackage{array}
\usepackage{longtable}
\usepackage{float}
\usepackage{amsmath}
\usepackage{caption}
% A opção `table` liga o colortbl, que é o que dá a zebra (\rowcolors).
\usepackage[table]{xcolor}
\usepackage{enumitem}
% Capa: a curva em TikZ e os ícones vetoriais do Font Awesome.
\usepackage{tikz}
\usepackage{fontawesome5}

% ── AS FONTES: serifa no corpo, Lato em tabela e figura ───────────────
% `lato` SEM a opção `default`: o corpo não muda e o comando \lato fica
% disponível. A `libertine` entra DEPOIS da `lato`, porque as duas mexem
% em \sfdefault e a última vence. A matemática acompanha a Libertine.
\usepackage{lato}
\usepackage{libertine}
\usepackage{libertinust1math}

% ── AS CORES: ColorBrewer Set1 e a escala tirada do #0A0310 ───────────
% Os mesmos hexadecimais do estilo_portfolio.py e do theme_portfolio.R.
\definecolor{azul}{HTML}{377EB8}       % destaque principal
\definecolor{vermelho}{HTML}{E41A1C}   % destaque secundário
\definecolor{verde}{HTML}{4DAF4A}
\definecolor{roxo}{HTML}{984EA3}
\definecolor{laranja}{HTML}{FF7F00}
\definecolor{marrom}{HTML}{A65628}
\definecolor{rosa}{HTML}{F781BF}
\definecolor{tinta}{HTML}{0A0310}      % o quase-preto do kit
\definecolor{tinta2}{HTML}{403A45}     % GRAY_70: texto secundário
\definecolor{tinta3}{HTML}{5D5961}     % GRAY_60: rodapé, nota
% Versões de TEXTO das cores claras, escurecidas até 4,5:1 contra o branco
% por cor_de_texto() do estilo_portfolio.py. Vermelho e roxo já passam.
\definecolor{azulTexto}{HTML}{3579B1}
\definecolor{verdeTexto}{HTML}{3C823B}

% ── A ZEBRA, aplicada a TODA tabela sem editar tabela nenhuma ─────────
% A alternância é ancorada no PRIMEIRO \midrule de cada tabular (o fim do
% cabeçalho), e não no número da linha: assim a primeira linha de dado
% sai sempre branca, com cabeçalho de uma ou de duas linhas.
%
% ⚠ A ordem {zebraclara}{white} é contraintuitiva e está certa: o
% contador sobe ANTES de a cor ser escolhida. Invertê-la inverte a zebra.
% ⚠ O reset vai dentro de \noalign, senão "Misplaced \noalign".
\usepackage{etoolbox}
\definecolor{zebraclara}{HTML}{F4F4F4}
\newif\ifzebraprimeiromidrule
\newcommand{\zebratabela}{%
  \lato
  \global\zebraprimeiromidruletrue
  \rowcolors{1}{zebraclara}{white}%
}
\AtBeginEnvironment{tabular}{\zebratabela}
\AtBeginEnvironment{longtable}{\zebratabela}
\pretocmd{\midrule}{%
  \noalign{%
    \ifzebraprimeiromidrule
      \global\zebraprimeiromidrulefalse
      \global\rownum=0\relax
    \fi}%
}{}{\typeout{ATENCAO: o patch de \string\midrule\space para a zebra FALHOU}}

% ── A BOLINHA DE ESTADO, para coluna sim/não de tabela ────────────────
% É reforço, e não substituição: o texto fica ao lado. Os três tons têm o
% mesmo contraste (2,3:1), calculado no projeto de origem. Glifo, e não
% TikZ: \tikz dentro de \caption derruba a compilação.
\definecolor{bolinhaSim}{HTML}{6DBD6B}
\definecolor{bolinhaNao}{HTML}{F29192}
\definecolor{bolinhaEmpate}{HTML}{B1AE28}
\DeclareRobustCommand{\bolinha}[1]{\textcolor{#1}{%
  \raisebox{-0.06ex}{\scalebox{2}{$\bullet$}}}}
\DeclareRobustCommand{\marcasim}{\bolinha{bolinhaSim}\,sim}
\DeclareRobustCommand{\marcanao}{\bolinha{bolinhaNao}\,não}
\DeclareRobustCommand{\marcaempate}{\bolinha{bolinhaEmpate}\,empate}

% ── A CAIXA DE DESTAQUE ───────────────────────────────────────────────
% Para o achado principal de uma seção, uma por seção no máximo. Fio azul
% à esquerda, fundo da zebra, título em Lato. Quebra entre páginas.
%   \begin{destaque}[Título opcional] texto \end{destaque}
\usepackage[most]{tcolorbox}
\newtcolorbox{destaque}[1][Em uma frase]{%
  enhanced, breakable,
  colback=zebraclara, colframe=azul,
  boxrule=0pt, leftrule=2.5pt, arc=0pt, outer arc=0pt,
  left=10pt, right=10pt, top=8pt, bottom=8pt,
  before skip=12pt, after skip=12pt,
  % ⚠ A cor do título é `coltitle`: um \color antes de \tcbtitle não vale,
  % porque o tcolorbox reaplica coltitle (branco por padrão) por dentro.
  title={#1}, detach title, coltitle=azulTexto,
  before upper={{\lato\bfseries\tcbtitle}\par\smallskip},
}

\usepackage{needspace}
\usepackage[hidelinks]{hyperref}
\setcounter{tocdepth}{2}
\hypersetup{bookmarksdepth=3, bookmarksnumbered=true, bookmarksopen=false}

\captionsetup{font=small, labelfont=bf}
\captionsetup[sub]{font=footnotesize}

\graphicspath{{figuras/}{exemplos/}}

% Paginação: altura natural, sem esticar o espaço vertical; sem órfãs nem
% viúvas; folga mínima de composição em vez de reescrever prosa.
\raggedbottom
\widowpenalty=10000
\clubpenalty=10000
\emergencystretch=1em

% ── \figwidth: a figura é mais larga que a prosa ──────────────────────
% A figura vive na medida de 1 cm de margem no papel (19 cm em A4), e não
% na mancha de 15 cm. Definida pelo papel, e não por número copiado.
\newlength{\figwidth}
\setlength{\figwidth}{\dimexpr\paperwidth-2cm\relax}

% \imagemplena[fração]{arquivo}: a imagem centrada sobre a mancha. A caixa
% de largura declarada impede o transbordo só para a direita e não conta
% como Overfull \hbox.
\newcommand{\imagemplena}[2][1]{%
  \makebox[\linewidth][c]{\includegraphics[width=#1\figwidth]{#2}}}

% tabelaplena: a tabela que não cabe na mancha ganha a medida da figura.
% Só para tabular (caixa não quebra entre páginas).
\newsavebox{\tabelacaixa}
\newenvironment{tabelaplena}
  {\begin{lrbox}{\tabelacaixa}}
  {\end{lrbox}\makebox[\linewidth][c]{\usebox{\tabelacaixa}}}

% \subfig[largura]{arquivo}{legenda}: um painel por linha, na medida da figura.
\newcommand{\subfig}[3][0.62\figwidth]{%
  \makebox[\linewidth][c]{%
  \begin{subfigure}[t]{#1}
    \centering
    \includegraphics[width=\linewidth]{#2}
    \caption{#3}
  \end{subfigure}}\par\medskip}
% \linhasde{arquivo}: as linhas de uma tabela, escritas por script.
% ⚠ Não use \input dentro de tabular: ele não é expansível, e o
% \bottomrule seguinte derruba a compilação com "Misplaced \noalign". O
% \@@input primitivo é expansível, e o espaço depois de #1 encerra o nome.
\makeatletter
\newcommand{\linhasde}[1]{\@@input #1 }
\makeatother
% ══════════════════════════════════════════════════════════════════════
% FIM DO PREÂMBULO DO KIT
% ══════════════════════════════════════════════════════════════════════

% ── AS FIGURAS DO PROJETO ─────────────────────────────────────────────
% O \graphicspath do kit (dentro do preâmbulo, que não se edita) aponta para
% figuras/ e exemplos/. Este vem depois e vence: compilando em relatorio/,
% as figuras escritas por scripts/05_figuras.py estão em ../figuras/.
\graphicspath{{../figuras/}}

% ── DADOS DA CAPA: preencha por relatório ─────────────────────────────
\newcommand{\RotuloCapa}{RELATÓRIO DE AVALIAÇÃO \textperiodcentered\ PROJETO DE PORTFÓLIO}
\newcommand{\TituloRelatorio}{Assistente do relatório Gestão Hospitalar no SUS-SP}
\newcommand{\SubtituloRelatorio}{Quanto um assistente com RAG acerta, medido com intervalo de confiança e comparação pareada}
% A segunda linha vem do registro: sobre as Rodadas reais, diz que os Veredictos foram revisados; sobre as
% sintéticas, que os números são inventados (scripts/04_analisar.py escolhe o texto).
\newcommand{\DelimitacaoRelatorio}{Relatório-fonte de \NumCorpusPaginas{} páginas e \NumGabaritoNPerguntas{} Perguntas \\ \NumRodadasOrigem}
\newcommand{\AutorNome}{João Eduardo Pastori Garcia}
\newcommand{\AutorPapel}{Estatístico responsável}
\newcommand{\AutorLinkedIn}{https://www.linkedin.com/in/joao-eduardo-garcia/}
\newcommand{\AutorGitHub}{https://github.com/JoaoEPGarcia}
\newcommand{\LocalCapa}{São Paulo -- SP}
% O mês e o ano da compilação, e não um ano digitado (GUIA_ESTILO.md, "A regra dos números").
\newcommand{\DataCapa}{\ifcase\month\or Janeiro\or Fevereiro\or Março\or Abril\or Maio\or Junho\or
  Julho\or Agosto\or Setembro\or Outubro\or Novembro\or Dezembro\fi\ de \number\year}

% ── O DESENHO DA CAPA ─────────────────────────────────────────────────
% A curva e os quatro traços da capa do kit, iguais aos do modelo, guardados
% aqui numa macro. Por quê: as coordenadas do TikZ são dígitos, e o
% conferidor de números (scripts/conferir_numeros.py) lê o corpo do .tex, do
% início do documento em diante, atrás de número digitado. Coordenada de
% desenho não é número impresso; no preâmbulo, ela fica fora da leitura sem
% que o conferidor precise de exceção, e o texto corrido continua conferido.
\newcommand{\DesenhoDaCapa}{%
  \begin{tikzpicture}
    \draw[line width=1.6pt, azul, smooth]
      plot coordinates {(0.2,0.35) (1.7,0.62) (3.2,0.48) (4.7,0.95)
                        (6.2,1.32) (7.7,1.12) (9.2,1.58) (10.7,1.42)};
    \foreach \x/\c in {0.2/verde, 2.9/vermelho, 5.6/azul, 8.3/roxo}
      \fill[\c] (\x,-0.28) rectangle ++(2.4,0.07);
  \end{tikzpicture}}

% ── OS NÚMEROS DO RELATÓRIO: gerados por script, nunca digitados ──────
% Cada \Num... vem de numeros.tex, escrito por scripts/06_gerar_documentos.py
% a partir de resultados/numeros.json. Ver GUIA_ESTILO.md, "A regra dos números".
% GERADO por scripts/06_gerar_documentos.py a partir de resultados/numeros.json.
% Não edite à mão: mude o script que calcula o número e rode o pipeline de novo.
\newcommand{\NumCorpusLegendas}{104}
\newcommand{\NumCorpusPaginas}{174}
\newcommand{\NumCorpusPalavras}{60.370}
\newcommand{\NumCorpusTabelas}{28}
\newcommand{\NumCorpusTrechos}{570}
\newcommand{\NumCorpusTrechosLegenda}{109}
\newcommand{\NumCorpusTrechosTabela}{71}
\newcommand{\NumCorpusTrechosTexto}{390}
\newcommand{\NumQuebraAlvo}{180}
\newcommand{\NumQuebraSobreposicao}{30}


% Frases que dependem de uma contagem do registro: o \ifcase escolhe a forma
% (singular, plural, nenhum) sem digitar número no texto.
\newcommand{\RevisaoMudou}{%
  \ifcase\NumRevisaoVeredictosMudados\relax
    a Revisão concordou com o Juiz em todos%
  \else
    a Revisão mudou \NumRevisaoVeredictosMudados{} deles%
  \fi}
\newcommand{\JuizDaRodada}{%
  \ifcase\NumParametrosNModelosJuiz\relax
  \or
    O Juiz foi o \texttt{\NumParametrosModeloJuiz}.%
  \else
    O Juiz usou \NumParametrosNModelosJuiz{} modelos em cadeia, nesta
    ordem: \NumParametrosModeloJuiz; quando a cota diária gratuita de um
    acabava, o seguinte assumia.%
  \fi}

\title{\TituloRelatorio}

\begin{document}

% ══════════════════════════════ CAPA ══════════════════════════════════
% Azul como destaque principal e vermelho como secundário. As quatro cores
% principais da paleta aparecem juntas em um só lugar: os traços ao pé da
% curva. Nenhum outro elemento da capa usa verde ou roxo.
\begin{titlepage}
  \centering
  \vspace*{0.5cm}

  {\color{tinta}\rule{\linewidth}{1.5pt}} \\[0.4cm]
  {\small\bfseries\color{tinta2} \RotuloCapa} \\[0.2cm]
  {\color{tinta}\rule{\linewidth}{0.5pt}}

  \vspace{1.6cm}

  % A curva ecoa a série temporal; ela e os quatro traços cobrem o mesmo vão.
  \DesenhoDaCapa

  \vspace{1.2cm}

  \begin{minipage}{0.94\textwidth}
    \centering
    {\Huge\bfseries\color{azul} \TituloRelatorio \par}
    \vspace{0.5cm}
    {\LARGE\bfseries\color{vermelho} \SubtituloRelatorio \par}
  \end{minipage}

  \vspace{1.1cm}

  % Um ícone por eixo do estudo: o Relatório-fonte, a busca, o modelo e a estatística.
  % ⚠ Use os nomes sem sufixo `Alt` (\faMapMarker, \faCalendarDay): o
  % fontawesome5 do MiKTeX não define os `...Alt`.
  {\Large
   \textcolor{azul}{\faDatabase}\hspace{1.1cm}
   \textcolor{vermelho}{\faSearch}\hspace{1.1cm}
   \textcolor{azul}{\faRobot}\hspace{1.1cm}
   \textcolor{vermelho}{\faCalculator}\par}

  \vspace{1.0cm}

  {\large\itshape\color{tinta2} \DelimitacaoRelatorio \par}

  \vspace{2.2cm}

  {\large\bfseries\color{azul} Autoria} \\[0.6cm]

  \begin{minipage}{0.85\textwidth}
    \centering
    {\large\bfseries \AutorNome} \\
    {\normalsize\color{tinta2} \AutorPapel} \\[0.15cm]
    {\footnotesize
     \href{\AutorLinkedIn}{\color{vermelho}\textbf{LinkedIn}} \ \textperiodcentered\ \
     \href{\AutorGitHub}{\color{vermelho}\textbf{GitHub}}}
  \end{minipage}

  \vfill

  {\color{tinta2}\rule{4cm}{0.5pt}} \\[0.3cm]
  {\small\color{tinta2} \LocalCapa \\ \bfseries \DataCapa}
\end{titlepage}

\tableofcontents
\clearpage

% O aviso de números inventados só existe sobre Rodadas sintéticas: sobre as
% reais, \NumAvisoSintetico é vazio e a caixa não aparece.
\ifdefempty{\NumAvisoSintetico}{}{%
  \begin{destaque}[Números inventados]
    \NumAvisoSintetico
  \end{destaque}}

% ══════════════════════════════════════════════════════════════════════
% CONTEÚDO DO PROJETO
% ══════════════════════════════════════════════════════════════════════
% As frases que dependem do sinal de um resultado (quem acerta mais, qual
% Classe de erro é a mais comum) não são escritas aqui: elas vêm de textos
% escolhidos por scripts/04_analisar.py a partir do resultado e registrados
% como macros (\NumComRagXSemRagConclusao, \NumComRagErroMaisFrequente...),
% ou dos títulos das figuras, escolhidos por scripts/05_figuras.py. O resto
% do texto continua verdadeiro com qualquer resultado.

\section{O problema}

Quase todo assistente com RAG (geração aumentada por recuperação) que se vê
em portfólio é uma demonstração sem medida: ele ``parece responder bem'', mas
ninguém sabe com que frequência acerta, nem se o acerto vem da busca no
texto consultado ou do que o modelo de linguagem já sabia antes de ler qualquer
coisa. Este projeto constrói um assistente que responde, em português,
perguntas sobre um relatório de pesquisa, o Relatório-fonte, e mede quanto ele acerta com as
ferramentas de um estatístico: intervalo de confiança, comparação pareada e
um desenho fixado antes de olhar a prova.

O texto consultado, chamado aqui de \textbf{Relatório-fonte}, é o
relatório ``Gestão Hospitalar no SUS-SP'', um PDF com publicação autorizada e
\NumCorpusPaginas{} páginas numeradas, além da capa. Um gestor de saúde
pergunta, por exemplo, quanto mudou um indicador numa rede de hospitais, e o
assistente responde em uma a três frases, cita a Página impressa de onde
tirou a resposta e mostra os Trechos em que se apoiou. Quando o
Relatório-fonte não trata do assunto, o comportamento certo é dizer isso, e
não inventar um número.

A pergunta do projeto tem três partes, e cada uma vira uma medida deste
Relatório de avaliação:

\begin{enumerate}[leftmargin=*, itemsep=2pt]
  \item \textbf{Quanto o assistente acerta?} O Acerto da resposta nas
    Perguntas respondíveis, com intervalo de Wilson, e não um número solto.
  \item \textbf{O acerto vem da busca?} A mesma prova é feita sem RAG, com o
    modelo respondendo só com o que já sabe, e as duas Condições são
    comparadas Pergunta a Pergunta.
  \item \textbf{Se o Relatório-fonte cabe inteiro no contexto do modelo, para que
    RAG?} Uma terceira Condição põe o Relatório-fonte inteiro no prompt, e o
    custo de cada Condição é medido em tokens de entrada e em latência.
\end{enumerate}

\begin{destaque}[O que este Relatório de avaliação entrega]
  O número de destaque é o Acerto da resposta com RAG nas
  \NumDestaqueNRespondiveis{} Perguntas respondíveis, com o intervalo de
  confiança de \NumIcNivel\%, ao lado do mesmo acerto sem RAG. Não se usa um
  ``acerto geral'' sobre todas as \NumGabaritoNPerguntas{} Perguntas: sem
  RAG, abster-se de tudo inflaria essa taxa, porque a Abstenção é a resposta
  certa nas Perguntas sem resposta. Cada número impresso aqui sai de um
  script e pode ser refeito a partir das Rodadas versionadas.
\end{destaque}

O texto usa um vocabulário fixo, definido no glossário do repositório. Os
termos com inicial maiúscula (Pergunta, Trecho, Condição, Rodada, Gabarito,
Revisão, Veredicto, Juiz, Classe de erro) têm sempre o mesmo sentido.

\section{O assistente}

O assistente é um pipeline de cinco passos. Todos rodam a partir de scripts
versionados, e só o último chama um modelo de linguagem.

\paragraph{Extração.} O \texttt{pdfplumber} lê o Relatório-fonte página a
página, com normalização Unicode (NFKC) para desfazer as ligaduras e com a
hifenização de fim de linha desfeita. Cada linha é classificada pelo layout
do LaTeX em que o Relatório-fonte foi escrito: corpo, título de capítulo,
título de seção ou legenda. As tabelas são detectadas à parte e
serializadas linha a linha, com cada célula acompanhada do nome da sua
coluna; uma tabela que continua na página seguinte herda a legenda. A
capa fica fora. A extração somou \NumCorpusPalavras{} palavras,
\NumCorpusTabelas{} blocos de tabela e \NumCorpusLegendas{} legendas de
figura.

\paragraph{Quebra em Trechos.} O texto vira Trechos, os pedaços que o
assistente pode recuperar. Nenhum Trecho atravessa a fronteira de página,
para que a citação aponte uma página só. Um parágrafo que cabe no alvo de
\NumQuebraAlvo{} palavras vai inteiro; um maior é dividido por frase, com
até \NumQuebraSobreposicao{} palavras repetidas entre Trechos vizinhos da
mesma página e seção. O título da seção abre cada Trecho, e cada tabela e
cada legenda formam Trechos próprios. São \NumCorpusTrechos{} Trechos:
\NumCorpusTrechosTexto{} de texto corrido, \NumCorpusTrechosTabela{} de
tabela e \NumCorpusTrechosLegenda{} de legenda.

\paragraph{Busca.} Cada Trecho vira um vetor do modelo
\texttt{intfloat/multilingual-e5-large}, rodado localmente pelo
\texttt{fastembed}, sem API paga. O índice é o FAISS com busca exata
(\texttt{IndexFlatIP} sobre vetores normalizados, ou seja, cosseno exato). Com
poucas centenas de Trechos, a busca exata custa quase nada e é
determinística: a mesma Pergunta recupera sempre os mesmos Trechos, o que
torna o Acerto da recuperação reprodutível. Na Condição com RAG, os
\NumParametrosK{} Trechos mais próximos da Pergunta vão ao modelo, cada um
marcado com a sua Página e o seu tipo (texto, tabela ou legenda).

\paragraph{Resposta.} O respondedor é o
\texttt{\NumParametrosModeloRespondedor}, do Google, chamado pela API do
Google AI Studio com saída estruturada por JSON Schema: a resposta, as
páginas citadas e um campo que marca a Abstenção. O prompt manda copiar os
números como estão no Relatório-fonte, no formato brasileiro, e não menciona
nenhum fato do Relatório-fonte, o que vazaria a resposta na Condição sem
RAG. Cada resposta registra o modelo que de fato respondeu, os tokens e a
latência.

\paragraph{API e Docker.} O assistente é servido por uma API FastAPI, com
\texttt{POST /perguntar} e \texttt{GET /saude}, empacotada em Docker. O
modelo de embeddings e o índice ficam em volumes, fora da imagem, e são
criados na primeira subida. O CI roda o lint, os testes e o build da imagem a
cada envio ao repositório.

\section{O desenho da avaliação}

\subsection{Perguntas, Gabarito e Revisão}

A prova tem \NumGabaritoNPerguntas{} Perguntas:
\NumGabaritoNRespondiveis{} respondíveis, divididas em
\NumGabaritoNTipos{} Tipos de pergunta conforme o lugar onde a resposta
está (número na prosa, número em tabela, número em legenda, conceito ou
método, e síntese entre páginas), e \NumGabaritoNSemResposta{} sem resposta
no Relatório-fonte, em que o certo é a Abstenção. As respondíveis foram
espalhadas pelos capítulos em proporção ao volume de texto.

As Perguntas e o Gabarito de cada uma (a resposta certa e as páginas que a
sustentam) foram \textbf{propostos com ajuda de um modelo de linguagem e
conferidos pelo autor} diretamente no Relatório-fonte, linha a linha. Essa conferência é
a Revisão: o script de avaliação recusa um Gabarito com qualquer linha sem
ela e lista as pendentes. Antes da Revisão, uma conferência mecânica compara
o Gabarito com o texto extraído (formato, páginas existentes, valores
presentes na página citada) para pegar erro grosseiro. Quando a mesma
resposta aparece em mais de uma página (o sumário executivo e a conclusão
repetem achados dos capítulos), o Gabarito lista todas, para que uma
citação certa não conte como errada.

\subsection{As três Condições}

As mesmas Perguntas passam pelo mesmo modelo em três Condições, que mudam
só o que o modelo recebe:

\begin{itemize}[leftmargin=*, itemsep=2pt]
  \item \textbf{com RAG}: os \NumParametrosK{} Trechos recuperados pela busca;
  \item \textbf{sem RAG}: nada do Relatório-fonte, só a Pergunta; mede o que o
    modelo já sabe, ou acha que sabe;
  \item \textbf{contexto completo}: o texto inteiro do Relatório-fonte, página
    por página, no prompt de sistema, com o cache de prompt da API.
\end{itemize}

A comparação principal é com RAG contra sem RAG: ela responde quanto do
acerto vem da busca. A comparação com o contexto completo responde se a
busca ainda se justifica quando o Relatório-fonte cabe na janela do modelo.

\subsection{As medidas}

Cada Pergunta, em cada Condição, recebe três medidas binárias, e cada taxa
sai com o intervalo de Wilson de \NumIcNivel\%, que se comporta bem com
amostras pequenas e taxas perto dos extremos:

\begin{itemize}[leftmargin=*, itemsep=2pt]
  \item \textbf{Acerto da resposta}, a medida principal: o conteúdo da
    resposta coincide com o Gabarito;
  \item \textbf{Acerto da citação}: alguma página citada está entre as
    páginas do Gabarito;
  \item \textbf{Acerto da recuperação}: alguma página do Gabarito está entre
    as dos Trechos recuperados; só existe na Condição com RAG, e separa
    problema de busca de problema de leitura.
\end{itemize}

As Perguntas respondíveis e as sem resposta são contadas em separado. Nas
sem resposta, a medida é a Abstenção correta. Cada resposta errada recebe
uma Classe de erro, derivada das medidas, sem julgamento novo, e as Classes
são mutuamente exclusivas: \textbf{resposta inventada} (respondeu a uma
Pergunta sem resposta), \textbf{citação errada} (a resposta está certa, a
página não), \textbf{abstenção indevida} (absteve-se numa respondível),
\textbf{falha de recuperação} (errou e a busca não trouxe nenhuma página do
Gabarito) e \textbf{falha de leitura} (errou com a página certa em mãos, ou
numa Condição que não recupera).

\subsection{O Veredicto: regra, Juiz e Revisão}

Os Tipos número na prosa, número em tabela e número em legenda têm resposta
numérica. Para essas Perguntas, o Veredicto sai
de uma regra determinística: todos os valores do Gabarito precisam aparecer,
exatamente, na resposta. A regra é estrita de propósito, porque o prompt
manda copiar os números como estão. Quando ela reprova, a resposta vai
também ao Juiz, para pegar o falso negativo da regra (uma diferença
expressa com o sinal trocado e a palavra ``para baixo'', por exemplo). As
respostas descritivas vão sempre ao Juiz, um modelo de linguagem que
compara o conteúdo com o Gabarito e ignora as páginas. \JuizDaRodada

O Juiz não tem a última palavra. Na Rodada principal, ele deu
\NumRevisaoVeredictosJuiz{} Veredictos, e o autor os conferiu numa folha;
a Revisão vence o Juiz. Foram \NumRevisaoVeredictosRevisados{} Veredictos
revisados, e \RevisaoMudou. A análise final se
recusa a publicar enquanto houver Veredicto do Juiz sem Revisão.

\subsection{Pré-registro}

O número de Trechos recuperados, o tamanho do Trecho, os modelos, o nível
de raciocínio e os prompts ficam num arquivo versionado de parâmetros. Eles
foram ajustados só com Perguntas de desenvolvimento, que nunca entram no
Gabarito, e congelados por commit antes da Rodada principal: ajustar o
sistema olhando a prova inflaria a taxa de acerto. O script de avaliação
recusa a Rodada principal enquanto os parâmetros não estiverem congelados e
commitados. A Rodada principal registrou o commit em que rodou,
\texttt{\NumParametrosCommit}.

\begin{destaque}[As regras de confiança viram código]
  \begin{itemize}[leftmargin=*, itemsep=1pt, topsep=2pt]
    \item A avaliação recusa um Gabarito com qualquer linha sem Revisão.
    \item A Rodada principal recusa parâmetros que não estejam congelados por
      commit.
    \item Todo Veredicto do Juiz passa pela Revisão do autor antes de virar
      resultado.
    \item Todo número impresso sai de um registro alimentado por script, e um
      conferidor falha diante de número digitado no README ou neste
      Relatório de avaliação.
  \end{itemize}
\end{destaque}

\section{Resultados com intervalos}

A Tabela~\ref{tab:condicoes} traz as três medidas nas três Condições, cada
uma com o intervalo de Wilson. A Rodada principal fez
\NumPrincipalNRespostas{} respostas (cada Pergunta em cada Condição). As
três primeiras colunas valem para as \NumDestaqueNRespondiveis{} Perguntas
respondíveis; a última, para as \NumGabaritoNSemResposta{} sem resposta.
O Acerto da recuperação só existe com RAG.

\begin{table}[H]
  \centering
  \caption{As medidas por Condição, na Rodada principal, com o intervalo de
    Wilson de \NumIcNivel\% entre parênteses. Acertos da resposta, da
    citação e da recuperação sobre as Perguntas respondíveis; Abstenção
    correta sobre as Perguntas sem resposta.}
  \label{tab:condicoes}
  \footnotesize
  \begin{tabelaplena}
  \begin{tabular}{@{}lllll@{}}
    \toprule
    Condição & Acerto da resposta & Acerto da citação & Acerto da recuperação & Abstenção correta \\
    \midrule
    \linhasde{tabelas/resultados_por_condicao.tex}
    \bottomrule
  \end{tabular}
  \end{tabelaplena}
\end{table}

\begin{destaque}[O número de destaque]
  Com RAG, o Acerto da resposta foi de \NumDestaqueAcertoComRag{} nas
  \NumDestaqueNRespondiveis{} Perguntas respondíveis (IC de \NumIcNivel\%:
  \NumDestaqueAcertoComRagInf{} a \NumDestaqueAcertoComRagSup), contra
  \NumDestaqueAcertoSemRag{} sem RAG. Pelo intervalo da diferença pareada
  (Seção~\ref{sec:pareada}), a leitura é: \NumComRagXSemRagConclusao.
\end{destaque}

A Figura~\ref{fig:acerto} mostra o mesmo Acerto da resposta, com a linha
fina marcando o intervalo. O título da figura é escolhido pelo script a
partir do resultado, e afirma o achado que os números sustentam.

\begin{figure}[H]
  \centering
  \imagemplena[0.86]{acerto_por_condicao}
  \caption{Acerto da resposta nas Perguntas respondíveis, por Condição, com
    o intervalo de Wilson de \NumIcNivel\%. A Condição com RAG sai na cor de
    destaque, e as outras em cinza.}
  \label{fig:acerto}
\end{figure}

Dois cuidados de leitura. Primeiro, intervalos que se sobrepõem não provam
que duas Condições empatam, nem intervalos separados bastam para
compará-las: as Condições responderam às mesmas Perguntas, e a comparação
certa é a pareada, da próxima seção. Segundo, o Acerto da citação mede só
se a página citada está entre as do Gabarito, e por isso pode ser alto
mesmo numa Condição que não recebe o texto, se o modelo souber ou
adivinhar a página.

\section{Comparação pareada}
\label{sec:pareada}

Como as três Condições responderam às mesmas Perguntas, cada comparação usa
os pares: para cada Pergunta respondível, as duas Condições acertam, só uma
acerta, ou as duas erram. Só os pares discordantes informam a diferença. O
teste é o de McNemar exato, sobre os pares discordantes, e o intervalo da
diferença pareada entre os acertos é o de Newcombe (método do escore, para
proporções pareadas), também de \NumIcNivel\%.

\begin{table}[H]
  \centering
  \caption{Comparação pareada do Acerto da resposta nas Perguntas
    respondíveis da Rodada principal. A diferença é com RAG menos a outra
    Condição, em pontos percentuais, com o intervalo de Newcombe de
    \NumIcNivel\%; o valor-p é o do teste exato de McNemar.}
  \label{tab:pareada}
  \footnotesize
  \begin{tabelaplena}
  \begin{tabular}{@{}lrrrrll@{}}
    \toprule
    & As duas & Só com RAG & Só a outra & As duas & Diferença & Valor-p \\
    Comparação & acertam & acerta & acerta & erram & (IC de \NumIcNivel\%) & (McNemar) \\
    \midrule
    \linhasde{tabelas/comparacao_pareada.tex}
    \bottomrule
  \end{tabular}
  \end{tabelaplena}
\end{table}

Na comparação principal, foram \NumComRagXSemRagNPares{} pares: em
\NumComRagXSemRagSoComRag{} deles só a Condição com RAG acertou, e em
\NumComRagXSemRagSoSemRag{} só a Condição sem RAG acertou. A diferença
pareada foi de \NumComRagXSemRagDiferenca{} (IC de \NumIcNivel\%:
\NumComRagXSemRagDiferencaInf{} a \NumComRagXSemRagDiferencaSup), com
valor-p de McNemar \NumComRagXSemRagPMcnemar. Contra o contexto
completo, os pares discordantes ficaram assim:
\NumComRagXContextoCompletoSoComRag{} a favor de com RAG e
\NumComRagXContextoCompletoSoContextoCompleto{} a favor do contexto
completo; a diferença foi de \NumComRagXContextoCompletoDiferenca{} (IC:
\NumComRagXContextoCompletoDiferencaInf{} a
\NumComRagXContextoCompletoDiferencaSup), com valor-p
\NumComRagXContextoCompletoPMcnemar.

\begin{destaque}[O que o intervalo pareado permite dizer]
  \begin{itemize}[leftmargin=*, itemsep=1pt, topsep=2pt]
    \item Com RAG contra sem RAG: \NumComRagXSemRagConclusao.
    \item Com RAG contra contexto completo: \NumComRagXContextoCompletoConclusao.
  \end{itemize}
  A leitura segue o intervalo da diferença: todo acima de zero, todo abaixo
  de zero, ou contendo o zero. Quando o intervalo contém o zero, o resultado
  é declarado como empate não resolvido, e não como prova de igualdade.
\end{destaque}

\section{Erros por Tipo de pergunta}

A Tabela~\ref{tab:erros} cruza as Classes de erro com o Tipo de pergunta,
na Condição com RAG. Em cada linha estão o número de Perguntas do Tipo, os
acertos e as respostas em cada Classe de erro. A resposta inventada só pode
acontecer nas Perguntas sem resposta, e as outras Classes, só nas
respondíveis; daí os traços. A citação errada conta respostas certas com a
página errada, e por isso uma linha pode somar mais erros que Perguntas
erradas.

\begin{table}[H]
  \centering
  \caption{Classes de erro por Tipo de pergunta, na Condição com RAG, Rodada
    principal. Na linha das Perguntas sem resposta, os acertos são as
    Abstenções corretas.}
  \label{tab:erros}
  \small
  \begin{tabelaplena}
  \begin{tabular}{@{}lrrrrrrr@{}}
    \toprule
    & & & Falha de & Falha de & Abstenção & Citação & Resposta \\
    Tipo de pergunta & $n$ & Acertos & recuperação & leitura & indevida & errada & inventada \\
    \midrule
    \linhasde{tabelas/erros_por_tipo.tex}
    \bottomrule
  \end{tabular}
  \end{tabelaplena}
\end{table}

As Classes separam as causas que pedem correções diferentes. A falha de
recuperação aponta para a busca: o Trecho certo não chegou ao modelo, e o
remédio seria mudar o número de Trechos, o tamanho do Trecho ou o modelo de
embeddings. A falha de leitura aponta para o modelo: o Trecho certo estava
lá, e a resposta saiu errada. A abstenção indevida e a resposta inventada
são as duas faces da calibração: abster-se demais ou de menos.

\begin{destaque}[A Classe de erro mais frequente com RAG]
  \begin{itemize}[leftmargin=*, itemsep=1pt, topsep=2pt]
    \item Classe mais frequente: \NumComRagErroMaisFrequente.
    \item Respostas nessa Classe (em cada uma, se houver empate):
      \NumComRagErroMaisFrequenteN.
    \item Respostas com alguma Classe de erro, nas respondíveis e nas sem
      resposta: \NumComRagNComErro.
  \end{itemize}
\end{destaque}

A Figura~\ref{fig:erros} mostra o mesmo cruzamento, um painel por Classe de
erro, com a Classe mais frequente em destaque.

\begin{figure}[H]
  \centering
  \imagemplena[0.92]{erros_por_tipo}
  \caption{Classes de erro por Tipo de pergunta, na Condição com RAG. Cada
    painel é uma Classe; as barras contam Perguntas. Entre parênteses, o
    número de Perguntas do Tipo e o total de erros da Classe.}
  \label{fig:erros}
\end{figure}

Com poucas Perguntas por Tipo, a diferença entre dois Tipos raramente sai do
ruído: a tabela serve para apontar onde olhar primeiro, e não para ordenar os
Tipos com confiança.

\section{Tokens e latência}

Sem fatura em dólar por Pergunta, o custo de cada Condição é medido em
tokens de entrada, que é o que cresce com o texto enviado ao modelo, e a
espera, em latência por Pergunta. Por construção, a Condição contexto
completo envia o texto inteiro do Relatório-fonte a cada Pergunta; o cache
de prompt da API reaproveita esse prefixo, e os tokens lidos do cache
somaram \NumContextoCompletoTokensCacheLeituraTotal{} na Rodada principal.

\begin{table}[H]
  \centering
  \caption{Custo por Condição na Rodada principal: tokens de entrada em
    média por Pergunta (inclusive os lidos do cache), latência mediana por
    Pergunta e respostas dadas por outro modelo que não o pedido (fallback
    do servidor).}
  \label{tab:custo}
  \small
  \begin{tabular}{@{}lrrr@{}}
    \toprule
    Condição & Tokens de entrada (média) & Latência mediana (s) & Fallbacks \\
    \midrule
    \linhasde{tabelas/custo_por_condicao.tex}
    \bottomrule
  \end{tabular}
\end{table}

\begin{destaque}[Se cabe no contexto, para que RAG?]
  Com RAG, cada Pergunta levou em média \NumComRagTokensEntradaMedio{}
  tokens de entrada, com latência mediana de \NumComRagLatenciaMedianaS{} s;
  com o contexto completo, \NumContextoCompletoTokensEntradaMedio{} tokens e
  \NumContextoCompletoLatenciaMedianaS{} s. Entre as duas, o intervalo
  pareado do Acerto da resposta diz: \NumComRagXContextoCompletoConclusao.
\end{destaque}

A Figura~\ref{fig:custo} mostra a distribuição por Pergunta, e não só a
média: cada caixa resume as Perguntas de uma Condição, com o acerto ao lado
do nome. O eixo dos tokens é logarítmico, porque as Condições diferem em
ordens de grandeza.

\begin{figure}[H]
  \centering
  \imagemplena[0.92]{custo_por_condicao}
  \caption{Tokens de entrada e latência por Pergunta, por Condição, na
    Rodada principal, com o Acerto da resposta ao lado de cada Condição. O
    traço na caixa é a mediana, e o losango, a média.}
  \label{fig:custo}
\end{figure}

Na Rodada principal inteira, as três Condições somaram
\NumPrincipalTokensEntradaTotal{} tokens de entrada e
\NumPrincipalTokensSaidaTotal{} de saída. Respostas dadas por um modelo
diferente do pedido (fallback do servidor): \NumPrincipalNFallback.

\section{Variabilidade entre Rodadas}

O projeto não fixa a temperatura do modelo, então duas Rodadas iguais podem
dar Veredictos diferentes. \NumRepeticoesLeitura{} O número de repetições
comparadas com a Rodada principal é \NumRepeticoesComparadas. As repetições
não passam pela Revisão do autor: a concordância compara os Veredictos
automáticos (da regra ou do Juiz) e, para a comparação ser justa, a Rodada
principal entra do mesmo jeito, sem a Revisão. A Tabela~\ref{tab:concordancia}
dá, para cada repetição e Condição, a fração das Perguntas em que o
Veredicto foi o mesmo.

\begin{table}[H]
  \centering
  \caption{Concordância dos Veredictos automáticos entre cada repetição e a
    Rodada principal, por Condição, sobre todas as Perguntas comparadas.
    Veredictos não revisados.}
  \label{tab:concordancia}
  \small
  \begin{tabular}{@{}llrr@{}}
    \toprule
    Repetição & Condição & Concordância & Perguntas \\
    \midrule
    \linhasde{tabelas/concordancia.tex}
    \bottomrule
  \end{tabular}
\end{table}

A concordância mede a estabilidade do sistema inteiro, respondedor e Juiz
juntos. Qualquer discordância entre Rodadas diz que parte do acerto
publicado depende do sorteio do modelo, e é por isso que o
número de destaque vem da Rodada revisada, e não da média das Rodadas.

\section{Limitações}

O que os dados deste Relatório de avaliação não autorizam afirmar:

\begin{itemize}[leftmargin=*, itemsep=3pt]
  \item \textbf{Só texto.} O assistente lê o texto, as tabelas e as legendas
    do Relatório-fonte, mas não os gráficos: um valor que só existe dentro
    da imagem de um gráfico não é extraído. As Perguntas que dependeriam
    deles ficaram fora do Gabarito, e por isso o acerto medido não vale
    para elas.
  \item \textbf{Amostra pequena.} São \NumGabaritoNRespondiveis{} Perguntas
    respondíveis e \NumGabaritoNSemResposta{} sem resposta. O intervalo do
    número de destaque vai de \NumDestaqueAcertoComRagInf{} a
    \NumDestaqueAcertoComRagSup, e o cruzamento por Tipo de pergunta tem
    poucas Perguntas por célula: diferenças entre Tipos são indícios, e não
    resultados.
  \item \textbf{Um documento.} Tudo foi medido sobre um único Relatório-fonte,
    em português, com o layout do Relatório-fonte, escrito em LaTeX. O acerto não se
    transfere para outro documento, outro layout ou outro idioma sem medir
    de novo.
  \item \textbf{Gabarito proposto com ajuda de um modelo.} As Perguntas e as
    respostas do Gabarito foram propostas com ajuda de um modelo de
    linguagem e conferidas pelo autor no Relatório-fonte. A Revisão corrige erros, mas
    não apaga o viés de escolha: Perguntas que um modelo acha naturais de
    propor podem ser mais fáceis para um modelo responder.
  \item \textbf{Veredictos revisados por uma pessoa.} Dos
    \NumRevisaoVeredictosJuiz{} Veredictos do Juiz na Rodada principal,
    \NumRevisaoVeredictosRevisados{} foram revisados pelo autor, e
    \RevisaoMudou. A Revisão é de um revisor só, sem medida
    de concordância entre revisores. \JuizDaRodada{} O respondedor foi o
    \texttt{\NumParametrosModeloRespondedor}.
  \item \textbf{Repetições não revisadas.} A concordância entre Rodadas usa
    Veredictos automáticos e mede a estabilidade, e não o acerto.
  \item \textbf{Um modelo.} Os números publicados são de um só modelo, o
    \texttt{\NumParametrosModeloRespondedor}. O provedor local (Ollama)
    existe no código, mas não gera número publicado.
  \item \textbf{Custo em tokens.} O custo é medido em tokens de entrada e em
    latência, e não em dinheiro; a latência depende da carga do servidor
    no momento da Rodada.
\end{itemize}

\clearpage
\appendix
\section{Como refazer os números}

As Rodadas ficam versionadas no repositório, com o commit e os parâmetros de
cada execução ao lado. Um comando refaz tudo a partir delas, sem chamar o
modelo nem gastar cota:

\begin{center}
  \texttt{uv run python scripts/reproduzir.py}
\end{center}

Ele refaz a análise, as figuras, as macros deste Relatório de avaliação, o
README e a lista de números canônicos, compila este PDF e roda o conferidor
de números, parando no primeiro passo que falhar. Cada número impresso aqui
tem uma linha em \texttt{NUMEROS\_CANONICOS.md}, com o script que o calcula
e os lugares onde aparece.

\end{document}
